% !TEX program = xelatex
\documentclass[aspectratio=169,11pt,handout]{beamer}
\usetheme{Madrid}
\usepackage{fontspec}
\setsansfont{DejaVu Sans}
\setmainfont{DejaVu Serif}
\usepackage{polyglossia}
\setdefaultlanguage{turkish}
\usepackage{amsmath,amssymb}
\usepackage{booktabs}
\usepackage{multicol}
\usepackage{xcolor}

% ---------- Renkler ----------
\definecolor{anamavi}{RGB}{0,45,98}
\definecolor{cozumyesil}{RGB}{0,110,50}
\definecolor{uyarikirmizi}{RGB}{180,20,20}
\setbeamercolor{structure}{fg=anamavi}
\setbeamercolor{frametitle}{bg=anamavi,fg=white}
\setbeamercolor{title}{bg=anamavi,fg=white}
\setbeamercolor{palette primary}{bg=anamavi,fg=white}
\setbeamercolor{palette secondary}{bg=anamavi!85!black,fg=white}
\setbeamercolor{palette tertiary}{bg=anamavi!70!black,fg=white}
\setbeamercolor{block title}{bg=anamavi,fg=white}
\setbeamercolor{block body}{bg=anamavi!7}
\setbeamercolor{block title example}{bg=cozumyesil,fg=white}
\setbeamercolor{block body example}{bg=cozumyesil!8}
\setbeamercolor{block title alerted}{bg=uyarikirmizi,fg=white}
\setbeamercolor{block body alerted}{bg=uyarikirmizi!7}
\setbeamertemplate{navigation symbols}{}
\setbeamertemplate{blocks}[rounded][shadow=false]

% ---------- Soru / Çözüm kalıbı ----------
% \soru{başlık}{soru metni}{çözüm metni}: önce soru, bir tıklamayla çözüm
\newcommand{\soru}[3]{%
\begin{frame}{#1}
  \begin{block}{Soru}\large #2\end{block}
  \pause
  \begin{exampleblock}{Çözüm}\large #3\end{exampleblock}
\end{frame}}
\newcommand{\ab}[1]{\overline{#1}}

\title[MMAT105 -- Hafta 1]{MMAT105 Genel Matematik}
\subtitle{Hafta 1: Temel Kavramlar, İşlem Önceliği, İşaret ve Tek--Çift, Ardışık Sayılar, Faktöriyel}
\author{}
\date{2026--2027 Güz Dönemi}

\begin{document}

% =====================================================
\begin{frame}
\titlepage
\end{frame}

% (iletişim çerçevesi kaldırıldı)

\begin{frame}{Dönem Planı}
\scriptsize
\begin{columns}[t]
\column{0.49\textwidth}
\begin{tabular}{@{}r p{0.86\linewidth}@{}}
\toprule
Hf & Konu\\\midrule
1 & Temel kavramlar, işlem önceliği, ardışık sayılar, faktöriyel\\
2 & Asal sayılar, asal çarpanlar, tam bölenler, bölme\\
3 & Bölünebilme kuralları, EBOB -- EKOK\\
4 & Rasyonel ve ondalık sayılar\\
5 & 1. derece bir bilinmeyenli denklemler\\
6 & 1. derece iki bilinmeyenli denklemler\\
7 & Basit eşitsizlikler, mutlak değer\\
8 & \textbf{Ara Sınav}\\\bottomrule
\end{tabular}
\column{0.49\textwidth}
\begin{tabular}{@{}r p{0.86\linewidth}@{}}
\toprule
Hf & Konu\\\midrule
9 & Üslü sayılar\\
10 & Köklü sayılar\\
11 & Çarpanlara ayırma ve özdeşlikler\\
12 & Oran -- orantı, ortalamalar\\
13 & Sayı -- kesir ve yaş problemleri\\
14 & Yüzde, kâr -- zarar, karışım, faiz; tekrar\\
15 & \textbf{Final Sınavı}\\\bottomrule
\end{tabular}
\end{columns}
\vspace{6pt}
\begin{block}{Dersin işleyişi}
Kısa kural \(\rightarrow\) çözümlü örnekler \(\rightarrow\) sınıfta siz çözün \(\rightarrow\) haftalık çalışma kağıdı
\end{block}
\end{frame}

\begin{frame}{Bugün}
\begin{columns}[t]
\column{0.32\textwidth}
\begin{block}{Ders 1}
\begin{enumerate}
\item Sayı kümeleri: \(\mathbb{N}\), \(\mathbb{Z}\)
\item Rakam ve sayı
\item Basamak değeri
\item Çözümleme
\item En büyük -- en küçük
\end{enumerate}
\end{block}
\column{0.32\textwidth}
\begin{block}{Ders 2}
\begin{enumerate}\setcounter{enumi}{5}
\item İşlem önceliği
\item İşaret kuralları
\item Tek -- çift sayılar
\end{enumerate}
\end{block}
\column{0.32\textwidth}
\begin{block}{Ders 3}
\begin{enumerate}\setcounter{enumi}{8}
\item Ardışık sayılar: terim sayısı, toplam
\item Özel toplamlar
\item Faktöriyel, sadeleştirme
\item Sondaki sıfırlar, birler basamağı
\end{enumerate}
\end{block}
\end{columns}
\end{frame}

\begin{frame}{Isınma: Kafadan Hesap}
\begin{columns}[t]
\column{0.5\textwidth}
\textbf{Hesaplayınız:}
\begin{enumerate}
\item \(7+(-5)\)
\item \(5-(-8)\)
\item \(3+4\cdot 2\)
\item \((-2)\cdot(-6)\)
\item \(20\div 4\cdot 5\)
\end{enumerate}
\column{0.4\textwidth}
\pause
\textbf{Cevaplar:}
\begin{enumerate}
\item \(2\)
\item \(13\)
\item \(11\) \quad(\(14\) değil!)
\item \(12\)
\item \(25\) \quad(\(1\) değil!)
\end{enumerate}
\end{columns}
\end{frame}

% =====================================================
\section{Ders 1: Sayı Kümeleri, Rakam ve Basamak}
\begin{frame}
\centering\Large\color{anamavi}\textbf{Ders 1}\\[6pt]\normalsize Sayı Kümeleri, Rakam ve Basamak
\end{frame}

\subsection{Sayı kümeleri}
\begin{frame}{Sayı Kümeleri}
\begin{block}{Tanımlar}
\begin{itemize}
\item Doğal sayılar: \(\mathbb{N}=\{0,1,2,3,\dots\}\)
\item Sayma sayıları: \(\mathbb{N}^{+}=\{1,2,3,\dots\}\)
\item Tam sayılar: \(\mathbb{Z}=\{\dots,-2,-1,0,1,2,\dots\}\)
\item \(\mathbb{Z}^{+}=\{1,2,3,\dots\}\), \quad \(\mathbb{Z}^{-}=\{\dots,-3,-2,-1\}\), \quad \(\mathbb{Z}=\mathbb{Z}^{-}\cup\{0\}\cup\mathbb{Z}^{+}\)
\item \(\mathbb{N}\subset\mathbb{Z}\)
\end{itemize}
\end{block}
\begin{alertblock}{Dikkat}
\(0\) bir doğal sayıdır; \(0\) ne pozitif ne negatiftir; \(0\) çift bir sayıdır.
\end{alertblock}
\end{frame}

\soru{Sayı Kümeleri -- Örnek 1}
{\(-3,\ 0,\ 7,\ 2{,}5,\ \dfrac{12}{4},\ -\dfrac{8}{2}\) sayılarından hangileri doğal sayı, hangileri tam sayıdır?}
{\(\dfrac{12}{4}=3\), \(-\dfrac{8}{2}=-4\) olduğundan:
\begin{itemize}
\item Doğal sayılar: \(0,\ 7,\ \dfrac{12}{4}\)
\item Tam sayılar: \(-3,\ 0,\ 7,\ \dfrac{12}{4},\ -\dfrac{8}{2}\) \quad (\(2{,}5\) tam sayı değildir)
\end{itemize}}

\soru{Sayı Kümeleri -- Örnek 2}
{En küçük doğal sayı, en küçük pozitif tam sayı ve en büyük negatif tam sayı toplamı kaçtır?}
{En küçük doğal sayı \(0\), en küçük pozitif tam sayı \(1\), en büyük negatif tam sayı \(-1\).
\[0+1+(-1)=0\]}

\soru{Sayı Kümeleri -- Örnek 3}
{\(-4\) ile \(3\) arasında kaç tane tam sayı vardır? Uçlar dahil edilirse kaç olur?}
{Arada: \(-3,-2,-1,0,1,2\) \(\Rightarrow\) \(6\) tane.\\[3pt]
Uçlar dahil: \(3-(-4)+1=8\) tane.\\[3pt]
\textbf{Genel kural:} \(a<b\) tam sayıları için \([a,b]\) aralığında \(b-a+1\) tam sayı vardır.}

\subsection{Rakam ve basamak}
\begin{frame}{Rakam, Sayı, Basamak}
\begin{block}{Tanımlar}
\begin{itemize}
\item \textbf{Rakam:} Sayıları yazmakta kullanılan semboller: \(\{0,1,2,\dots,9\}\) \quad (10 tane)
\item \textbf{Sayı:} Rakamların bir çokluk belirtecek biçimde yan yana yazılması.
\item \textbf{Basamak değeri:} Rakamın, bulunduğu basamağın değeriyle çarpımı.
\end{itemize}
\end{block}
\begin{block}{Çözümleme}
\[\ab{ab}=10a+b,\qquad \ab{abc}=100a+10b+c,\qquad \ab{abcd}=1000a+100b+10c+d\]
Burada \(a\neq 0\) (ilk rakam sıfır olamaz).
\end{block}
\end{frame}

\soru{Basamak Değeri -- Örnek 4}
{\(2403\) sayısının binler, yüzler, onlar ve birler basamaklarındaki rakamların basamak değerleri nedir?}
{\[2403=2\cdot 1000+4\cdot 100+0\cdot 10+3\cdot 1\]
Binler: \(2000\), \quad Yüzler: \(400\), \quad Onlar: \(0\), \quad Birler: \(3\)}

\soru{Basamak Değeri -- Örnek 5}
{\(3050\) sayısının rakamları toplamı ile basamak değerleri toplamı arasındaki fark kaçtır?}
{Rakamlar toplamı: \(3+0+5+0=8\)\\[3pt]
Basamak değerleri toplamı: \(3000+0+50+0=3050\) \quad (sayının kendisi!)\\[3pt]
Fark: \(3050-8=3042\)}

\soru{En Büyük -- En Küçük: Örnek 6}
{İki basamaklı en büyük doğal sayı ile iki basamaklı en küçük doğal sayının toplamı kaçtır?}
{En büyük: \(99\), \quad en küçük: \(10\)
\[99+10=109\]}

\soru{En Büyük -- En Küçük: Örnek 7}
{Rakamları birbirinden farklı üç basamaklı en büyük doğal sayı ile en küçük doğal sayının farkı kaçtır?}
{En büyük: \(987\). \quad En küçük: \(102\) \quad (ilk rakam \(0\) olamaz; \(1\), sonra \(0\), sonra \(2\)).
\[987-102=885\]}

\soru{En Büyük -- En Küçük: Örnek 8 (Tuzak)}
{Rakamları birbirinden farklı iki basamaklı en küçük \textbf{tam sayı} kaçtır?}
{Tam sayılarda negatifler de vardır; en küçük, mutlak değerce en büyük negatif sayıdır.\\[3pt]
Rakamları farklı iki basamaklı en büyük sayı \(98\) \(\Rightarrow\) cevap: \(-98\)\\[3pt]
\textit{Not:} Rakamları farklı olma şartı yoksa cevap \(-99\) olurdu.}

\soru{Sayma -- Örnek 9}
{Kaç tane iki basamaklı doğal sayı vardır? Bunlardan kaçının rakamları birbirinden farklıdır?}
{İki basamaklı sayılar: \(10,11,\dots,99\) \(\Rightarrow\) \(99-10+1=90\) tane.\\[3pt]
Rakamları farklı olanlar: Onlar basamağı \(9\) seçenek (\(1\)--\(9\)), birler basamağı kalan \(9\) seçenek:
\[9\cdot 9=81\]
(Kontrol: rakamları aynı olanlar \(11,22,\dots,99\): \(9\) tane; \(90-9=81\).)}

\soru{Çözümleme -- Örnek 10}
{\(\ab{ab}\) ve \(\ab{ba}\) iki basamaklı sayılardır. \(\ab{ab}+\ab{ba}=121\) ise \(a+b\) kaçtır?}
{\[\ab{ab}+\ab{ba}=(10a+b)+(10b+a)=11a+11b=11(a+b)\]
\[11(a+b)=121\ \Rightarrow\ a+b=11\]}

\soru{Çözümleme -- Örnek 11}
{\(\ab{ab}-\ab{ba}=54\) ise \(a-b\) kaçtır? Bu koşulu sağlayan en büyük \(\ab{ab}\) sayısı kaçtır?}
{\[(10a+b)-(10b+a)=9a-9b=9(a-b)=54\ \Rightarrow\ a-b=6\]
\(a\) en büyük olmalı: \(a=9\Rightarrow b=3\). \quad En büyük \(\ab{ab}=93\).\\[3pt]
(Kontrol: \(93-39=54\).)}

\soru{Çözümleme -- Örnek 12}
{\(\ab{abc}\) ve \(\ab{cba}\) üç basamaklı sayılardır. \(\ab{abc}-\ab{cba}=396\) ise \(a-c\) kaçtır?}
{\[(100a+10b+c)-(100c+10b+a)=99a-99c=99(a-c)\]
\[99(a-c)=396\ \Rightarrow\ a-c=4\]
Ortanca rakam \(b\) sonucu etkilemez.}

\soru{Çözümleme -- Örnek 13}
{İki basamaklı \(\ab{ab}\) sayısı, rakamları toplamının \(4\) katına eşittir. Bu koşulu sağlayan kaç sayı vardır?}
{\[10a+b=4(a+b)\ \Rightarrow\ 6a=3b\ \Rightarrow\ b=2a\]
\(a=1,2,3,4\) \(\Rightarrow\) \(12,\ 24,\ 36,\ 48\) \quad (\(a=5\) için \(b=10\) rakam değil)\\[3pt]
Cevap: \(4\) tane.}

\soru{Basamak Değişimi -- Örnek 14}
{Üç basamaklı \(\ab{abc}\) sayısında \(a\) rakamı \(2\) artırılır, \(c\) rakamı \(3\) azaltılırsa sayı nasıl değişir?}
{\(a\) yüzler basamağında: \(2\cdot 100=200\) artış.\\
\(c\) birler basamağında: \(3\cdot 1=3\) azalış.
\[200-3=197\ \text{artar.}\]}

\begin{frame}{Ders 1 -- Sınıfta Siz Çözün}
\begin{columns}[t]
\column{0.62\textwidth}
\begin{enumerate}
\item Rakamları farklı dört basamaklı en küçük doğal sayı kaçtır? \(\star\)
\item \(\ab{ab}+\ab{ba}=88\) ise \(a+b=?\) \(\star\)
\item \(\ab{ab}-\ab{ba}=45\) ise \(a-b=?\)
\item \(\ab{4a7}\) sayısında \(a\) yerine \(a+3\) yazılırsa sayı kaç artar? \(\star\)
\item \(-6\) ile \(5\) arasında (uçlar hariç) kaç tam sayı vardır?
\end{enumerate}
\column{0.34\textwidth}
\pause
\textbf{Cevaplar:}
\begin{enumerate}
\item \(1023\)
\item \(8\)
\item \(5\)
\item \(30\)
\item \(10\)
\end{enumerate}
\end{columns}
\vspace{4pt}{\footnotesize \(\star\) işaretli soruların adım adım çözümü sonraki slaytlardadır.}
\end{frame}

% ---- Ders 1 alıştırma çözümleri ----
\soru{Ders 1 -- Alıştırma Çözümü (Soru 1)}
{Rakamları birbirinden farklı dört basamaklı en küçük doğal sayı kaçtır?}
{Basamakları soldan sağa, en küçük uygun rakamla doldururuz:\\[3pt]
Binler: \(0\) olamaz \(\Rightarrow\) \(1\). \quad Yüzler: kullanılmamış en küçük rakam \(0\).\\[3pt]
Onlar: \(2\). \quad Birler: \(3\).
\[\text{Cevap: } 1023\]}

\soru{Ders 1 -- Alıştırma Çözümü (Soru 2)}
{\(\ab{ab}\) ve \(\ab{ba}\) iki basamaklı sayılardır. \(\ab{ab}+\ab{ba}=88\) ise \(a+b\) kaçtır?}
{Çözümleyelim:
\[\ab{ab}+\ab{ba}=(10a+b)+(10b+a)=11a+11b=11(a+b)\]
\[11(a+b)=88\ \Rightarrow\ a+b=8\]
(Örneğin \(35+53=88\).)}

\soru{Ders 1 -- Alıştırma Çözümü (Soru 4)}
{\(\ab{4a7}\) üç basamaklı sayısında \(a\) rakamı yerine \(a+3\) yazılırsa sayı kaç artar?}
{\(a\) rakamı onlar basamağındadır; basamak değeri \(10\cdot a\).\\[3pt]
Rakam \(3\) artınca basamak değeri \(10\cdot 3=30\) artar.
\[\big(400+10(a+3)+7\big)-\big(400+10a+7\big)=30\]
(\(a+3\) rakam olmalı: \(a\le 6\). Örneğin \(427\to 457\).)}

% =====================================================
\section{Ders 2: İşlem Önceliği, İşaret, Tek--Çift}
\begin{frame}
\centering\Large\color{anamavi}\textbf{Ders 2}\\[6pt]\normalsize İşlem Önceliği, İşaret Kuralları, Tek -- Çift Sayılar
\end{frame}

\subsection{İşlem önceliği}
\begin{frame}{İşlem Önceliği}
\begin{block}{Sıra}
\begin{enumerate}
\item Parantez: önce en içteki \(\ (\ )\ \rightarrow\ [\ ]\ \rightarrow\ \{\ \}\)
\item Üs alma
\item Çarpma ve bölme \quad (aynı öncelik, \textbf{soldan sağa})
\item Toplama ve çıkarma \quad (aynı öncelik, \textbf{soldan sağa})
\end{enumerate}
\end{block}
\begin{alertblock}{En sık yapılan hatalar}
\begin{itemize}
\item \(20\div 4\cdot 5=5\cdot 5=25\) \quad (\(20\div 20=1\) değil)
\item \(-3^2=-9\), \quad ama \quad \((-3)^2=9\)
\item Kesir çizgisi parantez gibidir: pay ve payda ayrı ayrı hesaplanır.
\end{itemize}
\end{alertblock}
\end{frame}

\soru{İşlem Önceliği -- Örnek 15}
{\(18-12\div 3\cdot 2+5\) işleminin sonucu kaçtır?}
{Önce bölme ve çarpma, soldan sağa: \(12\div 3=4\), \ \(4\cdot 2=8\)
\[18-8+5=15\]}

\soru{İşlem Önceliği -- Örnek 16}
{\(48\div 4\div 2\) ile \(48\div(4\div 2)\) işlemlerinin sonuçları nedir?}
{\[48\div 4\div 2=12\div 2=6\]
\[48\div(4\div 2)=48\div 2=24\]
Parantez sırayı değiştirir; bölme işleminde birleşme özelliği yoktur.}

\soru{İşlem Önceliği -- Örnek 17}
{\(5-3\cdot(2-6)+2^3\) işleminin sonucu kaçtır?}
{Parantez: \(2-6=-4\). \quad Üs: \(2^3=8\).
\[5-3\cdot(-4)+8=5+12+8=25\]}

\soru{İşlem Önceliği -- Örnek 18}
{\(3+2\cdot\big[10-(4-1)^2\big]\) işleminin sonucu kaçtır?}
{\[(4-1)^2=3^2=9,\qquad 10-9=1\]
\[3+2\cdot 1=5\]}

\soru{İşlem Önceliği -- Örnek 19}
{\(\big\{8-\big[6-(4-2)\big]\big\}\cdot 3\) işleminin sonucu kaçtır?}
{En içten dışa:
\[(4-2)=2\ \rightarrow\ [6-2]=4\ \rightarrow\ \{8-4\}=4\ \rightarrow\ 4\cdot 3=12\]}

\soru{İşlem Önceliği -- Örnek 20}
{\(\dfrac{12-3\cdot 2}{2\cdot 3-3}+2^3\cdot 3-4^2\div 8\) işleminin sonucu kaçtır?}
{Kesir: \(\dfrac{12-6}{6-3}=\dfrac{6}{3}=2\)\\[3pt]
\(2^3\cdot 3=24\), \quad \(4^2\div 8=16\div 8=2\)
\[2+24-2=24\]}

\subsection{İşaret kuralları}
\begin{frame}{İşaret Kuralları}
\begin{columns}[t]
\column{0.48\textwidth}
\begin{block}{Çarpma ve bölme}
\[(+)\cdot(+)=+\qquad (-)\cdot(-)=+\]
\[(+)\cdot(-)=-\qquad (-)\cdot(+)=-\]
Bölmede de aynı kural geçerlidir.
\end{block}
\column{0.48\textwidth}
\begin{block}{Kuvvetler (\(a>0\))}
\[(-a)^{2k}=a^{2k}>0\]
\[(-a)^{2k+1}=-a^{2k+1}<0\]
\[(-1)^{\text{çift}}=1,\qquad (-1)^{\text{tek}}=-1\]
\end{block}
\end{columns}
\begin{alertblock}{Çıkarma}
\(a-(-b)=a+b\): \ Negatif bir sayıyı çıkarmak, pozitifini eklemektir.
\end{alertblock}
\end{frame}

\soru{İşaret -- Örnek 21}
{\(-7+(-5)-(-12)\) ve \((-2)\cdot(-3)\cdot(-1)\cdot 4\) işlemlerinin sonuçları nedir?}
{\[-7-5+12=0\]
Üç negatif çarpan (tek sayıda) \(\Rightarrow\) sonuç negatif:
\[(-2)\cdot(-3)\cdot(-1)\cdot 4=-24\]}

\soru{İşaret -- Örnek 22}
{\((-3)^2-2^3\) ve \((-2)^4-(-3)^3\) işlemlerinin sonuçları nedir?}
{\[(-3)^2-2^3=9-8=1\]
\[(-2)^4-(-3)^3=16-(-27)=16+27=43\]}

\soru{İşaret -- Örnek 23}
{\(-2^2+(-2)^2\) ve \((-1)^{2026}+(-1)^{2027}-1^{2028}\) işlemlerinin sonuçları nedir?}
{\[-2^2+(-2)^2=-4+4=0\]
\[(-1)^{2026}+(-1)^{2027}-1^{2028}=1+(-1)-1=-1\]
Dikkat: \(1^{n}=1\) her \(n\) için; \(-1^{2028}=-(1^{2028})=-1\).}

\soru{İşaret -- Örnek 24}
{\(a<0<b\) olduğuna göre \(a\cdot b\), \ \(b-a\), \ \(a-b\), \ \(a^2\cdot b\) ifadelerinin işaretleri nedir?}
{\begin{itemize}
\item \(a\cdot b\): \((-)\cdot(+)\ \Rightarrow\ \) negatif
\item \(b-a=b+|a|\ \Rightarrow\ \) pozitif
\item \(a-b=-(b-a)\ \Rightarrow\ \) negatif
\item \(a^2>0\), \(b>0\ \Rightarrow\ a^2\cdot b\) pozitif
\end{itemize}}

\soru{İşaret -- Örnek 25}
{\(a\cdot b<0\), \quad \(b\cdot c>0\), \quad \(c<0\) olduğuna göre \(a\) ve \(b\)'nin işaretleri nedir?}
{\(b\cdot c>0\) ve \(c<0\ \Rightarrow\ b<0\)\\[3pt]
\(a\cdot b<0\) ve \(b<0\ \Rightarrow\ a>0\)\\[3pt]
Sonuç: \(a>0\), \ \(b<0\), \ \(c<0\)}

\soru{İşaret -- Örnek 26}
{\(a<0\) ve \(a^3\cdot b>0\) olduğuna göre \(b\)'nin işareti nedir?}
{\(a<0\Rightarrow a^3<0\) \ (tek kuvvet işareti korur).\\[3pt]
\(a^3\cdot b>0\) ve \(a^3<0\ \Rightarrow\ b<0\)}

\subsection{Tek ve çift sayılar}
\begin{frame}{Tek ve Çift Sayılar}
\begin{block}{Tanım}
Çift: \(2k\), \quad Tek: \(2k+1\) \qquad (\(k\in\mathbb{Z}\)). \quad \(0\) çifttir; negatif sayılar da tek/çift olabilir.
\end{block}
\begin{columns}[t]
\column{0.48\textwidth}
\begin{block}{Toplama / çıkarma}
\begin{tabular}{c|cc}
\(\pm\) & T & Ç\\\hline
T & Ç & T\\
Ç & T & Ç
\end{tabular}
\end{block}
\column{0.48\textwidth}
\begin{block}{Çarpma}
\begin{tabular}{c|cc}
\(\cdot\) & T & Ç\\\hline
T & T & Ç\\
Ç & Ç & Ç
\end{tabular}
\end{block}
\end{columns}
\vspace{4pt}
Kısa yol: Çarpanlardan en az biri çift ise çarpım çifttir. \(n\) tek ise \(n^k\) de tektir.
\end{frame}

\soru{Tek -- Çift: Örnek 27}
{\(x\) tek, \(y\) çift tam sayılardır. \(x+y\), \ \(x\cdot y\), \ \(x^2+y\), \ \(3x+2y\) ifadeleri tek mi çift mi?}
{\begin{itemize}
\item \(x+y\): T \(+\) Ç \(=\) \textbf{tek}
\item \(x\cdot y\): T \(\cdot\) Ç \(=\) \textbf{çift}
\item \(x^2+y\): T \(+\) Ç \(=\) \textbf{tek}
\item \(3x+2y\): T \(+\) Ç \(=\) \textbf{tek}
\end{itemize}}

\soru{Tek -- Çift: Örnek 28}
{\(x,y\in\mathbb{Z}\) ve \(3x+2y\) tek sayı ise \(x\) tek midir, çift midir? \(y\) hakkında ne söylenebilir?}
{\(2y\) her zaman çifttir. \(3x+\text{Ç}=\text{T}\ \Rightarrow\ 3x\) tek \(\Rightarrow\ x\) tek.\\[3pt]
\(y\) hakkında bilgi yoktur: \(y\) tek de olabilir, çift de.}

\soru{Tek -- Çift: Örnek 29}
{\(a,b,c\) tam sayılardır. \(a+b\) tek, \(b+c\) çift ise \(a+c\) tek midir, çift midir?}
{\(a+b\) tek \(\Rightarrow\) \(a\) ile \(b\) farklı türde.\\
\(b+c\) çift \(\Rightarrow\) \(b\) ile \(c\) aynı türde.\\
O halde \(a\) ile \(c\) farklı türde \(\Rightarrow\ a+c\) \textbf{tek}.\\[3pt]
Kontrol: \(a=1,\ b=2,\ c=4\): \(a+c=5\).}

\soru{Tek -- Çift: Örnek 30}
{\(a\cdot b\) çift ve \(a\) tek ise \(b\) için ne söylenebilir? \(a\cdot b\) tek ise \(a\) ve \(b\) için?}
{\(a\) tek ve çarpım çift \(\Rightarrow\) \(b\) \textbf{çift} olmak zorunda.\\[3pt]
Çarpım tek \(\Rightarrow\) çarpanların \textbf{hepsi tek}: \(a\) ve \(b\) tek.}

\begin{frame}[fragile]{Bilgisayar Köşesi: Python da aynı kuralları uygular}
\begin{columns}[t]
\column{0.5\textwidth}
\begin{verbatim}
>>> 20 / 4 * 5
25.0
>>> -3**2
-9
>>> (-3)**2
9
>>> 2**3**2
512
\end{verbatim}
\column{0.46\textwidth}
\begin{verbatim}
>>> 2403 % 10    # birler
3
>>> 2403 // 10 % 10  # onlar
0
>>> 7 % 2        # tek -> 1
1
\end{verbatim}
\end{columns}
\small \(2^{3^2}=2^{9}=512\): Üsler sağdan sola hesaplanır. \ \verb|%| kalanı, \verb|//| tam bölümü verir.
\end{frame}

\begin{frame}{Ders 2 -- Sınıfta Siz Çözün}
\begin{columns}[t]
\column{0.62\textwidth}
\begin{enumerate}
\item \(24\div 4\cdot 3-2^3=?\) \(\star\)
\item \(-3^2-(-3)^2+(-1)^{101}=?\) \(\star\)
\item \(a>0>b\) ise \(a\cdot b^2\) ve \(b-a\) işaretleri?
\item \(x\) tam sayı ve çift ise \(5x+7\) tek mi çift mi?
\item \(a\cdot b\cdot c<0\), \(a<0\), \(c>0\) ise \(b\)'nin işareti? \(\star\)
\end{enumerate}
\column{0.34\textwidth}
\pause
\textbf{Cevaplar:}
\begin{enumerate}
\item \(10\)
\item \(-19\)
\item pozitif; negatif
\item tek
\item \(b>0\)
\end{enumerate}
\end{columns}
\vspace{4pt}{\footnotesize \(\star\) işaretli soruların adım adım çözümü sonraki slaytlardadır.}
\end{frame}

% ---- Ders 2 alıştırma çözümleri ----
\soru{Ders 2 -- Alıştırma Çözümü (Soru 1)}
{\(24\div 4\cdot 3-2^3\) işleminin sonucu kaçtır?}
{Önce üs: \(2^3=8\).\\[3pt]
Çarpma ve bölme soldan sağa: \(24\div 4=6\), \quad \(6\cdot 3=18\).
\[18-8=10\]
(\(24\div 12=2\) yapmak hatalıdır.)}

\soru{Ders 2 -- Alıştırma Çözümü (Soru 2)}
{\(-3^2-(-3)^2+(-1)^{101}\) işleminin sonucu kaçtır?}
{\(-3^2=-(3^2)=-9\) \quad (üs yalnız \(3\)'e aittir)\\[3pt]
\((-3)^2=9\), \quad \((-1)^{101}=-1\) \ (tek kuvvet)
\[-9-9+(-1)=-19\]}

\soru{Ders 2 -- Alıştırma Çözümü (Soru 5)}
{\(a\cdot b\cdot c<0\), \ \(a<0\) ve \(c>0\) ise \(b\)'nin işareti nedir?}
{\(a\cdot c\): \((-)\cdot(+)\) \(\Rightarrow\) \(a\cdot c<0\).\\[3pt]
\(a\cdot b\cdot c=(a\cdot c)\cdot b<0\) ve \(a\cdot c<0\) \(\Rightarrow\) \(b\) pozitif olmalı.
\[b>0\]
(Kontrol: \(a=-1,\ b=2,\ c=3\): \(abc=-6<0\).)}

% =====================================================
\section{Ders 3: Ardışık Sayılar ve Faktöriyel}
\begin{frame}
\centering\Large\color{anamavi}\textbf{Ders 3}\\[6pt]\normalsize Ardışık Sayılar ve Faktöriyel
\end{frame}

\subsection{Ardışık sayılar}
\begin{frame}{Ardışık Sayılar: Terim Sayısı ve Toplam}
\begin{block}{Artışı sabit olan sayılar}
\[\text{Terim sayısı}=\frac{\text{Son terim}-\text{İlk terim}}{\text{Artış}}+1\]
\[\text{Toplam}=\text{Terim sayısı}\cdot\frac{\text{İlk terim}+\text{Son terim}}{2}\]
\end{block}
\begin{block}{Gösterim (\(x\in\mathbb{Z}\))}
Ardışık tam sayılar: \(x,\ x+1,\ x+2,\dots\) \qquad Ardışık çift / tek sayılar: \(x,\ x+2,\ x+4,\dots\)
\end{block}
\begin{alertblock}{Sık yapılan hata}
Terim sayısında ``\(+1\)'' unutulmamalı: \(3,5,7,\dots,21\) dizisinde \(\frac{21-3}{2}=9\) değil, \(9+1=10\) terim vardır.
\end{alertblock}
\end{frame}

\begin{frame}{Özel Toplamlar ve Ortanca Terim}
\begin{block}{Özel toplamlar (\(n\in\mathbb{N}^{+}\))}
\[1+2+3+\dots+n=\frac{n(n+1)}{2}\]
\[2+4+6+\dots+2n=n(n+1)\qquad\qquad 1+3+5+\dots+(2n-1)=n^2\]
\end{block}
\begin{block}{Ortanca terim}
Artışı sabit, \textbf{tek sayıda} terimin toplamı \(=\) terim sayısı \(\times\) ortanca terim.\\
Ortanca terim \(=\dfrac{\text{İlk terim}+\text{Son terim}}{2}\).
\end{block}
\end{frame}

\soru{Ardışık Sayılar -- Örnek 31}
{\(5+8+11+\dots+62\) toplamında kaç terim vardır? Toplam kaçtır?}
{Artış \(3\):
\[n=\frac{62-5}{3}+1=19+1=20\]
\[\text{Toplam}=20\cdot\frac{5+62}{2}=10\cdot 67=670\]}

\soru{Ardışık Sayılar -- Örnek 32}
{\(1+2+\dots+30\) toplamı kaçtır? \(11+12+\dots+30\) toplamı kaçtır?}
{\[1+2+\dots+30=\frac{30\cdot 31}{2}=465\]
\(11+\dots+30\): baştaki \(1+\dots+10=\dfrac{10\cdot 11}{2}=55\) çıkarılır:
\[465-55=410\]
(Kontrol: \(20\) terim, \(20\cdot\frac{11+30}{2}=410\).)}

\soru{Ardışık Sayılar -- Örnek 33}
{\(2+4+6+\dots+50\) ve \(1+3+5+\dots+49\) toplamları kaçtır?}
{Çiftler: \(2n=50\Rightarrow n=25\):
\[2+4+\dots+50=25\cdot 26=650\]
Tekler: \(2n-1=49\Rightarrow n=25\):
\[1+3+\dots+49=25^2=625\]}

\soru{Ardışık Sayılar -- Örnek 34}
{Ardışık \(5\) tam sayının toplamı \(85\) ise en büyüğü kaçtır?}
{Terim sayısı tek (\(5\)) \(\Rightarrow\) toplam \(=5\,\cdot\) ortanca terim.
\[\text{Ortanca}=\frac{85}{5}=17\]
Sayılar: \(15,16,17,18,19\) \(\Rightarrow\) en büyüğü \(19\).}

\soru{Ardışık Sayılar -- Örnek 35}
{Ardışık \(3\) çift sayının toplamı \(78\) ise en büyüğü kaçtır?}
{Sayılar \(x,\ x+2,\ x+4\):
\[3x+6=78\ \Rightarrow\ 3x=72\ \Rightarrow\ x=24\]
Sayılar: \(24,26,28\) \(\Rightarrow\) en büyüğü \(28\).\\[3pt]
(Kısa yol: ortanca \(78/3=26\).)}

\soru{Ardışık Sayılar -- Örnek 36}
{\(40\) ile \(120\) arasında (uçlar hariç) kaç tek sayı vardır? Bu sayıların toplamı kaçtır?}
{İlk tek sayı \(41\), son tek sayı \(119\), artış \(2\):
\[n=\frac{119-41}{2}+1=39+1=40\]
\[\text{Toplam}=40\cdot\frac{41+119}{2}=40\cdot 80=3200\]}

\soru{Ardışık Sayılar -- Örnek 37}
{Bir kitabın sayfaları \(1\)'den \(120\)'ye kadar numaralandırılmıştır. Bu iş için toplam kaç rakam kullanılır?}
{Basamak sayısına göre gruplayalım:
\begin{itemize}
\item \(1\)--\(9\): \(9\) sayı \(\times 1\) rakam \(=9\)
\item \(10\)--\(99\): \(90\) sayı \(\times 2\) rakam \(=180\)
\item \(100\)--\(120\): \(120-100+1=21\) sayı \(\times 3\) rakam \(=63\)
\end{itemize}
Toplam: \(9+180+63=252\) rakam.}

\subsection{Faktöriyel}
\begin{frame}{Faktöriyel}
\begin{block}{Tanım}
\[n!=1\cdot 2\cdot 3\cdots n\quad(n\in\mathbb{N}^{+}),\qquad 0!=1,\qquad n!=n\cdot(n-1)!\]
\centering
\begin{tabular}{cccccccc}
\toprule
\(1!\) & \(2!\) & \(3!\) & \(4!\) & \(5!\) & \(6!\) & \(7!\) & \(8!\)\\\midrule
\(1\) & \(2\) & \(6\) & \(24\) & \(120\) & \(720\) & \(5040\) & \(40320\)\\\bottomrule
\end{tabular}
\end{block}
\begin{block}{Sadeleştirme}
Büyük faktöriyeli küçüğe kadar açın: \(\dfrac{9!}{7!}=\dfrac{9\cdot 8\cdot 7!}{7!}=72\). \quad Toplamda küçük faktöriyeli dışarı alın: \(7!+6!=6!\,(7+1)=8\cdot 6!\).
\end{block}
\end{frame}

\begin{frame}{Faktöriyel: Sondaki Sıfırlar ve Birler Basamağı}
\begin{block}{Sondaki sıfır sayısı}
Her sıfır bir \(2\cdot 5\) çarpımından gelir; \(n!\) içinde \(2\) çarpanı \(5\)'ten çoktur. Bu yüzden \(5\) çarpanları sayılır:
\[\left\lfloor\frac{n}{5}\right\rfloor+\left\lfloor\frac{n}{25}\right\rfloor+\left\lfloor\frac{n}{125}\right\rfloor+\cdots\qquad(\lfloor x\rfloor:\ x\text{'in tam kısmı})\]
\end{block}
\begin{block}{Birler basamağı}
\(5!=120\) olduğundan \(n\ge 5\) için \(n!\)'in birler basamağı \(0\)'dır. Toplamlarda yalnız \(1!,2!,3!,4!\) terimlerine bakılır.
\end{block}
\begin{alertblock}{Sık yapılan hatalar}
\((a+b)!\neq a!+b!\), \quad \((2n)!\neq 2\cdot n!\), \quad \(\dfrac{6!}{3!}\neq 2!\) \ (\(\dfrac{6!}{3!}=120\)).
\end{alertblock}
\end{frame}

\soru{Faktöriyel -- Örnek 38}
{\(\dfrac{8!}{6!}\) ve \(\dfrac{7!}{4!\cdot 3!}\) işlemlerinin sonuçları nedir?}
{\[\frac{8!}{6!}=\frac{8\cdot 7\cdot 6!}{6!}=56\]
\[\frac{7!}{4!\cdot 3!}=\frac{7\cdot 6\cdot 5\cdot 4!}{4!\cdot 6}=\frac{210}{6}=35\]}

\soru{Faktöriyel -- Örnek 39}
{\(\dfrac{8!-7!}{6!}\) işleminin sonucu kaçtır?}
{Küçük faktöriyel \(7!\) ortak çarpan:
\[8!-7!=7!\,(8-1)=7\cdot 7!\]
\[\frac{7\cdot 7!}{6!}=\frac{7\cdot 7\cdot 6!}{6!}=49\]}

\soru{Faktöriyel -- Örnek 40}
{\((n+1)!=20\cdot(n-1)!\) olduğuna göre \(n\) kaçtır?}
{\((n+1)!=(n+1)\cdot n\cdot(n-1)!\) yazılırsa:
\[(n+1)\cdot n\cdot(n-1)!=20\cdot(n-1)!\ \Rightarrow\ n(n+1)=20=4\cdot 5\]
\[n=4\]
(Kontrol: \(5!=120=20\cdot 3!\).)}

\soru{Faktöriyel -- Örnek 41}
{\(100!\) sayısının sonunda kaç sıfır vardır?}
{\(5\) çarpanlarını sayalım:
\[\left\lfloor\frac{100}{5}\right\rfloor+\left\lfloor\frac{100}{25}\right\rfloor=20+4=24\]
(\(125>100\) olduğundan sonraki terim \(0\).)\\[3pt]
Cevap: \(24\) sıfır.}

\soru{Faktöriyel -- Örnek 42}
{\(n!\) sayısının sonunda tam \(6\) sıfır varsa \(n\) en az kaçtır?}
{\(n=20\): \(4\) sıfır; \ \(n=24\): \(4+0=4\) sıfır.\\[3pt]
\(n=25\): \(\left\lfloor\frac{25}{5}\right\rfloor+\left\lfloor\frac{25}{25}\right\rfloor=5+1=6\) sıfır.\\[3pt]
\(25\) iki tane \(5\) çarpanı getirdiği için sayı \(4\)'ten \(6\)'ya atlar (\(5\) sıfırlı faktöriyel yoktur).
\[n_{\min}=25\]}

\soru{Faktöriyel -- Örnek 43}
{\(1!+3!+5!+7!+\dots+21!\) toplamının birler basamağındaki rakam kaçtır? Toplam tek midir, çift midir?}
{\(5!\) ve sonrası birler basamağı \(0\) olan sayılardır.\\[3pt]
Önemli terimler: \(1!+3!=1+6=7\) \(\Rightarrow\) birler basamağı \(7\).\\[3pt]
Birler basamağı tek olduğundan toplam \textbf{tektir}.}

\begin{frame}{Ders 3 -- Sınıfta Siz Çözün}
\begin{columns}[t]
\column{0.62\textwidth}
\begin{enumerate}
\item \(4+9+14+\dots+99\) toplamı kaçtır? \(\star\)
\item Ardışık \(4\) çift sayının toplamı \(92\) ise en küçüğü kaçtır? \(\star\)
\item \(1+2+3+\dots+20=?\)
\item \(\dfrac{12!}{10!\cdot 2!}=?\)
\item \(75!\) sayısının sonunda kaç sıfır vardır? \(\star\)
\end{enumerate}
\column{0.34\textwidth}
\pause
\textbf{Cevaplar:}
\begin{enumerate}
\item \(1030\)
\item \(20\)
\item \(210\)
\item \(66\)
\item \(18\)
\end{enumerate}
\end{columns}
\vspace{4pt}{\footnotesize \(\star\) işaretli soruların adım adım çözümü sonraki slaytlardadır.}
\end{frame}

% ---- Ders 3 alıştırma çözümleri ----
\soru{Ders 3 -- Alıştırma Çözümü (Soru 1)}
{\(4+9+14+\dots+99\) toplamı kaçtır?}
{Artış \(5\). Terim sayısı:
\[n=\frac{99-4}{5}+1=19+1=20\]
\[\text{Toplam}=20\cdot\frac{4+99}{2}=10\cdot 103=1030\]}

\soru{Ders 3 -- Alıştırma Çözümü (Soru 2)}
{Ardışık \(4\) çift sayının toplamı \(92\) ise en küçüğü kaçtır?}
{Sayılar \(x,\ x+2,\ x+4,\ x+6\):
\[4x+12=92\ \Rightarrow\ 4x=80\ \Rightarrow\ x=20\]
Sayılar: \(20,22,24,26\) \quad (toplam \(92\) \checkmark).}

\soru{Ders 3 -- Alıştırma Çözümü (Soru 5)}
{\(75!\) sayısının sonunda kaç sıfır vardır?}
{\(5\) çarpanlarını sayalım:
\[\left\lfloor\frac{75}{5}\right\rfloor=15,\qquad \left\lfloor\frac{75}{25}\right\rfloor=3,\qquad \left\lfloor\frac{75}{125}\right\rfloor=0\]
\[15+3=18\ \text{sıfır}\]
(\(25,50,75\) ikişer tane \(5\) çarpanı içerir.)}

\begin{frame}{Özet}
\begin{columns}[t]
\column{0.32\textwidth}
\begin{block}{Ders 1}
\begin{itemize}
\item \(\mathbb{N}=\{0,1,2,\dots\}\), \(\mathbb{Z}\); \(0\) çift ve işaretsiz
\item \([a,b]\) aralığında \(b-a+1\) tam sayı
\item \(\ab{ab}\pm\ab{ba}\): \(11(a+b)\), \(9(a-b)\)
\item \(\ab{abc}-\ab{cba}=99(a-c)\)
\end{itemize}
\end{block}
\column{0.32\textwidth}
\begin{block}{Ders 2}
\begin{itemize}
\item Parantez \(\to\) üs \(\to\) \(\times\div\) \(\to\) \(+-\), soldan sağa
\item \(-a^2\neq(-a)^2\)
\item Tek sayıda negatif çarpan \(\Rightarrow\) negatif
\item Çift çarpan varsa çarpım çift
\end{itemize}
\end{block}
\column{0.32\textwidth}
\begin{block}{Ders 3}
\begin{itemize}
\item Terim sayısı: \(\dfrac{\text{son}-\text{ilk}}{\text{artış}}+1\)
\item \(\dfrac{n(n+1)}{2}\), \ \(n(n+1)\), \ \(n^2\)
\item \(n!=n\,(n-1)!\), \ \(0!=1\)
\item Sondaki sıfır: \(\lfloor\frac{n}{5}\rfloor+\lfloor\frac{n}{25}\rfloor+\cdots\)
\end{itemize}
\end{block}
\end{columns}
\vspace{6pt}
\begin{alertblock}{Gelecek hafta}
Asal sayılar, asal çarpanlara ayırma, tam bölenler, bölme
\end{alertblock}
\end{frame}

\begin{frame}
\centering\Large İyi çalışmalar
\end{frame}

\end{document}
